% c06 ------ 附录 A：黎曼几何中的旋量；附录 B：低维例子
% 源文件 QuantumWorkTry.tex 第 2480--2918 行（arXiv:1508.04715v2）
% 本文件仅为正文片段：公式、编号与引用按原稿逐字保留，逐段全译并附审校批注。

\appendix
\section{黎曼几何中的旋量（Spinors In Riemannian Geometry）}\label{dirram}

\begin{keytakeaways}
\begin{itemize}[leftmargin=1.6em,itemsep=0.2em,topsep=0.2em]
  \item \textbf{物理图像}：弯曲流形上的费米子不是普通函数，而是旋量丛的截面；能否整体定义旋量，取决于流形的定向性以及 spin、pin 结构的拓扑障碍。定向反转对称又通过反线性配对约束狄拉克谱，从而决定配分函数是实、复还是带模 1 相位。
  \item \textbf{数学工具}：用 SO(n) 与 O(n) 的双重覆盖 Spin(n)、pin+ 与 pin- 两支来拼接局部标架，过渡函数的符号模糊性由一阶 Z2 上同调分类；偶维时用手征算子构造自伴狄拉克算子，奇维时 pin+ 结构只能以互补对的形式出现。
  \item \textbf{核心结论}：spin 结构之间相差一阶 Z2 上同调的扭转；pin+ 与 pin- 结构的障碍类不同，前者为第二 Stiefel--Whitney 类，后者还须加上第一类的平方，故后者自四维起才出现障碍。偶维 pin+ 流形上存在自伴狄拉克算子；奇维则只能定义两个互补结构之间的指标，且该指标必为偶数。
\end{itemize}
\end{keytakeaways}

\subsection{可定向流形上的旋量（Spinors On Orientable Manifolds）}\label{spinor}

为在 $n$ 维黎曼流形 $W$ 上\footnote{在我们的应用中，$W$ 通常是 $D$ 维时空 $X$，或是其 $d=D-1$ 维边界 $Y=\partial X$。}引入旋量与狄拉克方程，方便的做法是先选取一个“标架”（vielbein），即一组正交归一的切向量基 $e_a$，$a=1,\dots,n$。记 $e_a^i$，$i=1,\dots,n$，为 $e_a$ 在 $W$ 上任意局部坐标系 $x^i=\{x^1,\dots,x^n\}$ 中的分量。$e_a$ 正交归一，意即：若 $g_{i j}$ 是 $W$ 的度规张量，则
\begin{equation}\label{orth} e_a^i e_b^jg_{ij}=\delta_{ab}.\end{equation}
等价地，
\begin{equation}\label{north} \delta^{ab}e_a^ie_b^j =g^{ij}. \end{equation}
$e_a^i$ 的指标用 $\delta_{ab}$ 与 $g_{ij}$ 升降，特别地 $e_i^a$ 是 $e_a^i$ 的逆矩阵。

% [Tutor Note: 旋量丛的局部平凡化本质上是在小开集上把 SO(n) 标架丛平凡化，再逐点提升到 Spin(n)；$e_a^i$ 与 $e_i^a$ 互为逆矩阵、指标分别用 $\delta_{ab}$ 与 $g_{ij}$ 升降，正是为了在弯曲流形上模拟平直空间的 $\{\g_a,\g_b\}=2\delta_{ab}$。原文把“标架”与“局部平凡化”混用，须记住真正的整体对象是过渡函数及其提升，而不是某个固定的 $e_a$；同一流形上不同标架给出的局部旋量描述只相差旋量表示中的逐点变换。]

“平直空间”的 $\\gamma$ 矩阵是 $2^{[n/2]}\times 2^{[n/2]}$ 矩阵 $\g_a, \,a=1,\dots,n$，满足
\begin{equation}\label{worth}\{\g_a,\g_b\}=2\delta_{ab}.\end{equation}
相应的“弯曲空间” $\\gamma$ 矩阵为 $\g_i=e_i^a\g_a$，满足
\begin{equation}\label{zorth}\{\g_i,\g_j\}=2g_{ij}.\end{equation}

% [Tutor Note: 欧几里得号差下通常取 $\g_a$ 厄米且 $(\g_a)^2=1$，由此得到 $\{\g_a,\g_b\}=2\delta_{ab}$；弯曲情形的 $\g_i=e_i^a\g_a$ 是 $\g_a$ 的实线性组合，仍厄米。狄拉克算子定义中显式的 $\i$ 因子不可省：紧致流形上协变导数 $D_i$ 反厄米，$\i$ 恰好抵消其反厄米性而使 $\D$ 自伴。不同文献对 $\g$ 取厄米或反厄米、对 $\{\g_i,\g_j\}$ 取 $+2g_{ij}$ 或 $-2g_{ij}$ 的约定不同，会使 $\i$ 的位置乃至符号随维度改变；对照公式前务必先对齐这一约定。]

在 $W$ 的给定点 $p\in W$ 处，切空间上作用的是保持定向的旋转群 $\SO(n)$，而 $W$ 上旋量对应的群是 $\SO(n)$ 的双重覆盖，称为 $\Spin(n)$。其生成元为
\begin{equation}\label{torth}\Sigma_{ab}=\frac{1}{4}[\g_a,\g_b]. \end{equation}
旋量场 $\Psi$ 的协变导数定义为
\begin{equation}\label{porth}\frac{D}{D x^i}\Psi=\left(\frac{\partial}{\partial x^i}+\omega_i^{ab}\Sigma_{ab}\right)\Psi,\end{equation}
其中 $\omega_i^{ab}$ 是 $W$ 切丛上的 Levi-Civita 联络。于是可以定义厄米的狄拉克算子：
\begin{equation}\label{dirac}\D=\i\sum_{i=1}^n\g^iD_i. \end{equation}

% [Tutor Note: $\Sigma_{ab}$ 的归一化 $\frac14[\g_a,\g_b]$ 使 $\Spin(n)$ 生成元在旋量表示中满足与 $\SO(n)$ 相同的李代数关系；协变导数中 $\omega_i^{ab}\Sigma_{ab}$ 的符号约定决定 $\D$ 是否自伴。原文取 $\D=\i\sum_i\g^iD_i$，其中 $\g^i=e^i_a\g^a$。若把 $\D$ 写成 $\i\slashed{D}$，注意 $\slashed{D}=\sum_i\g^iD_i$ 与 $\D$ 只差因子 $\i$，两种写法的自伴性与谱符号极易混淆。]

这一构造依赖于标架 $\{e_a\}$ 的选取，而标架一般只能局部地选取。要定义整体的狄拉克算子，须把局部描述拼合起来。一种朴素的做法是用小开集 $W_\x$，$\x=1,\dots,s$，覆盖 $W$，并在每个 $W_\x$ 上选取标架 $e_{a\,\x}$，$a=1,\dots,n$。在 $W_\x\cap W_\y$ 上，由于 $e_{a\,\x}$ 与 $e_{a\,\y}$ 都是 $W$ 切丛的正交归一基，它们由一个正交变换 $M_{\x \y}$ 相联系：$e_\x=M_{\x \y} e_\y$，更详细地写作
\begin{equation}\label{toroth}e_{a\,\x}=\sum_b M_a{}^b{}_{\x \y} e_{b\,\y}. \end{equation}
（这里 $M_a{}^b{}_{\x \y}$ 是 $n\times n$ 正交矩阵 $M_{\x \y}$ 的矩阵元。）若 $W$ 可定向，如本节所设，则可选取局部标架 $e_\x$ 使得 $M_{\x \y}$ 取值于 $\SO(n)$（而非 $\O(n)$）。过渡矩阵 $M_{\x \y}$ 满足相容性条件：在三重交叠 $W_\x\cap W_\y\cap W_\z$ 上有
\begin{equation}\label{horth} M_{\x \y}M_{\y \z}M_{\z \x}=1.\end{equation}

% [Tutor Note: 原文源第 2525 行把 vielbein 写作 vierbein，二者是同一德语词的不同拼写（“四维标架”），并非另一个概念；译文统一处理为“标架”。]

一旦在开集 $W_\x$ 中选取了局部标架，便可在此开集内定义旋量场 $\Psi_\x$（其分量为 $\Psi_{\alpha \,\x}$，$\alpha=1,\dots, 2^{[n/2]}$）。$\Psi_\x$ 的协变导数与狄拉克方程如式（\ref{orth}）和式（\ref{dirac}）所示。在交叠 $W_\x\cap W_\y$ 上，我们通过
\begin{equation}\label{compare}\Psi_\x =\h M_{\x \y}\Psi_\y,\end{equation}
来比较 $\Psi_\x$ 与 $\Psi_\y$，其中 $\h M_{\x \y}$ 是把 $M_{\x \y}$ 提升到旋量表示所得的 $2^{[n/2]}\times 2^{[n/2]}$ 矩阵。于是 $\h M_{\x \y}$ 是 $\Spin(n)$ 的元素，即与自旋相伴的 $\SO(n)$ 的双重覆盖。由于给定的 $\SO(n)$ 元素作用在旋量上的符号并非唯一确定，从 $M_{\x \y}$ 到 $\h M_{\x \y}$ 的提升只确定到相差一个符号。为相容起见，符号的选取必须使三重交叠满足
\begin{equation}\label{ompare}\h M_{\x \y}\h M_{\y \z}\h M_{\z \x}=1. \end{equation}

% [Tutor Note: 提升 $\h M_{\x\y}$ 的符号不由 $\SO(n)$ 元素唯一决定（$\Spin(n)\to\SO(n)$ 是双重覆盖），所以“自旋结构”的实质是三重交叠上符号的相容选取，而非某个局部矩阵。原文用 $(-1)^{c_{\x\y}}$ 刻画这一符号模糊性并证明 $c$ 的规范等价类给出 $H^1(W,\Z_2)$；这正是旋量丛区别于普通向量丛之处：旋量丛并非由切丛唯一决定，而是由切丛再加一个 $H^1(W,\Z_2)$ 扭转决定。]

满足此条件的符号选取称为 $W$ 上的一个自旋结构；一旦选定自旋结构，便可把局部描述拼接起来，在 $W$ 上整体地研究费米子与狄拉克方程。在维数 $\geq 4$ 时，给定的 $W$ 可能不容许自旋结构，因为可能无法满足式（\ref{ompare}）；此时 $W$ 上不能定义旋量（一个例子是 $W=\Bbb{CP}^2$）。这种情形在本文中不起主要作用。我们只需把注意力限于可定义旋量的时空。

对我们的目的而言更重要的是：给定的 $W$ 可能容许多个自旋结构。若确实找到了与式（\ref{compare}）相容的符号选取，则另一个选取 $\h M'_{\x \y} =\h M_{\x \y} (-1)^{c_{\x \y}}$（其中 $c_{\x \y}$ 是取值为 $\Z_2$ 的函数）满足同样的条件，只要对所有的 $\x,\y,\z$ 有
\be\label{yero}(-1)^{c_{\x \y}}(-1)^{c_{\y \z}}(-1)^{c_{\z \x}}=1\ee
我们希望给 $c$ 施加如下的规范等价：$c_{\x\y}\cong c_{\x\y}+d_\x-d_\y$，其中 $d_\x$ 是任意 $\Z_2$ 值函数；因为在该变换下，式（\ref{compare}）中由 $c_{\x\y}$ 改变所引起的 $\h M_{\x\y}$ 的变化，可以通过把局部旋量重新定义为 $\Psi_\x\to (-1)^{d_\x}\Psi_\x$ 来吸收。给定条件（\ref{yero}）与上述规范不变性，$c_{\x\y}$ 定义了群 $H^1(W,\Z_2)$ 中的一个元素。$W$ 上任意两个自旋结构相差一个由该群中元素所作的“扭转”。

一个初等的例子是取 $W$ 为带平坦度规的 $n$ 维环面 $T^n$。切向量在 $W$ 上的平行移动完全是平凡的：切向量上的 Levi-Civita 联络为 0。然而在 $W$ 上定义旋量时，我们可以自由地规定：旋量沿非可缩回路 $\ell\subset T^n$ 平行移动一周后变号。由于有 $n$ 个独立的回路需要考虑，在 $T^n$ 上定义旋量时有 $2^n$ 种符号选择，它们就是 $T^n$ 上的各个自旋结构。群 $H^1(T^n,\Z_2)$ 同构于 $\Z_2^n$，并对这些符号选择加以标号。（对 $W=T^n$，但并非一般情形，存在一个特殊的自旋结构，即旋量平行移动时不改变符号的“平凡”自旋结构。因此在这个例子中，自旋结构的集合可以典范地等同于群 $H^1(W,\Z_2)$。）

赋予自旋结构选取的流形 $W$ 称为自旋流形。（可定向性已包含在自旋流形的概念中，因为在出发点我们把切丛的过渡函数取为取值于 $\SO(n)$。）若 $W$ 是自旋流形，则 $M_{\x\y}$ 作为 $\Spin(n)$ 旋量表示中的 $2^{[n/2]}\times 2^{[n/2]}$ 矩阵，是向量丛 $\S\to W$ 的过渡函数。$W$ 上的费米子场是该丛的一个截面。我们有时把 $\S$ 称为自旋结构。若 $\S$ 是一个自旋结构，则任何其他自旋结构都是 $\S\otimes \rho$，其中 $\rho\to W$ 是二阶实线丛（在上述构造中由过渡函数 $(-1)^{c_{\x\y}}$ 定义）。

对偶数维 $n$，克利福德代数的不可约表示在同构意义下唯一。这可以通过把 $n$ 个 $\\gamma$ 矩阵写成 $n/2$ 对产生与湮灭算符来证明，而后者具有唯一的不可约表示。但对偶数 $n$，克利福德代数的不可约表示作为 $\Spin(n)$ 的表示并非不可约；它分解为正手征与负手征的两个 $\Spin(n)$ 表示之和。对奇数 $n$，情况不同。在同构意义下，$\Spin(n)$ 只有一个旋量表示，但克利福德代数有两个不等价的表示。这两个表示相差 $\g_i\to -\g_i$：后者保持克利福德代数的反对易关系，不影响 $\Spin(n)$ 生成元 $\Sigma_{ij}=\frac{1}{4}[\g_i,\g_j]$，却给出克利福德代数的一个不等价表示。这两个表示的差别在于乘积 $\g_1\g_2\cdots \g_n$ 的符号------该乘积与克利福德代数对易，因而在克利福德代数的不可约表示中是一个 $c$ 数（$\pm 1$ 或 $\pm \i$，视 $n$ 而定）。

例如在 3 维，克利福德代数的两个表示在局部欧氏标架下对应于 $\g_i=\pm \sigma_i$，从而 $\\gamma$ 矩阵满足
\be\label{zoffo}\g_i\g_j=g_{ij}\pm\i\epsilon_{ijk}\g^k, \ee
其中 $\epsilon_{ijk}$ 是 Levi-Civita 张量。在此式中，$\pm\epsilon_{ijk}$ 可视为 $W$ 的一个定向选取。于是克利福德代数的两个表示与 $W$ 的两种可能定向相关联。狄拉克算子 $\D=\i\sum_k\g^k D_k$ 在 $\g_k\to -\g_k$ 下为奇，因此狄拉克算子的符号依赖于 $W$ 的定向。若改变该算子的符号（从而改变其所有本征值的符号），狄拉克算子的 $\eta$ 不变量变号；这就是 $\eta$ 不变量在宇称（即定向反转）下为奇的原因。

% [Tutor Note: 奇维时 $\g_i\to-\g_i$ 不改变 $\Spin(n)$ 生成元，却使 $\D$ 变号；原文由此把克利福德代数的两个表示与两种定向等同起来。$\eta$ 不变量随 $\D\to-\D$ 变号，因为全部本征值变号，这正是 $\eta$ 在定向反转（宇称）下为奇的根源。这里也埋下了后文的伏笔：奇维不可定向流形上 $\P$ 与 $\P'$ 的模 2 指标由 $\eta$ 的模 2 约化给出，构成 3 维自旋拓扑量子场论分类的几何基础。]

\subsection{不可定向流形上的费米子（Fermions On An Unorientable Manifold）}\label{unor}

% [Tutor Note: $\Pin^\pm(n)$ 是 $\O(n)$ 的两个不同双重覆盖，其差别在于空间反射 $\sR$ 在旋量上的平方为 $+1$ 还是 $-1$（$\Pin^+$ 取 $+1$）。存在性障碍类的标准表述是：$\Pin^+$ 结构要求 $w_2=0$，$\Pin^-$ 结构要求 $w_2+w_1^2=0$（$w_i$ 为切丛的 Stiefel--Whitney 类，均取值于模 2 上同调）。故 $\RP^2$（$w_2=w_1^2\ne0$）没有 $\Pin^+$ 结构而有 $\Pin^-$ 结构；$\RP^4$（$w_2=0$ 而 $w_1^2\ne0$）有 $\Pin^+$ 结构而无 $\Pin^-$ 结构------这正是原文“$\pin^-$ 的障碍在四维才开始出现”一句的含义，也说明两种结构的存在性判断不能互换。]

现在考虑不可定向的情形。在不可定向的 $n$ 维流形 $W$ 上，$W$ 切丛的过渡矩阵 $M_{\x\y}$ 取值于 $\O(n)$ 而非 $\SO(n)$，相应地，作用在旋量上的矩阵 $\h M_{xy}$ 也必须取值于 $\O(n)$ 的某个双重覆盖。此双重覆盖有两种选择\footnote{这里我们考虑马约拉纳费米子，即只与引力耦合的费米子，这对拓扑超导体在物理上是适当的。在存在电磁作用的 $\U(1)$ 规范对称性时，可以把空间反射与一个规范变换结合起来，定义一个使 $\sR^2$ 以任意相位作用在电子上的对称性。在数学上，这与从 spin、pin 结构推广到 $\mathrm{spin}_c$、$\mathrm{pin}_c$ 结构的可能性有关。这一推广对分析拓扑物态有用，但本文中不会用到。}，取决于作用在旋量上的空间反射 $\sR$ 满足 $\sR^2=1$ 还是 $\sR^2=-1$。对大多数物理应用，适当的选择是 $\sR^2=1$。由 $\sR^2=1$ 产生的 $\Spin(n)$ 的双重覆盖称为 $\Pin^+(n)$。（若取 $\sR^2=-1$，则得到 $\Spin(n)$ 的一个称为 $\Pin^-(n)$ 的双重覆盖。）

% [Tutor Note: 原文此处说 $\Pin^\pm(n)$ 是“$\Spin(n)$ 的双重覆盖”，严格说不确：它们是 $\O(n)$ 的双重覆盖，$\Spin(n)$ 只是其单位分支（见本节开头及源第 2596 行的说法）。同一段还混用记号 $\h M_{xy}$ 与后文的 $\h M_{\x\y}$，应视为同一对象的两种写法。另需注意，$n=1$ 时 $\Pin^+(1)\cong\Z_2\times\Z_2$ 的覆盖可裂，故任意实线丛都有 $\pin^+$ 结构，而 $\Pin^-(1)\cong\Z_4$ 才有非平凡障碍------这是判断两族结构差别的一个有用检验。]

欧几里得号差下的反射对称性可以延拓为洛伦兹号差下的反射对称性或时间反演对称性，取决于该反射是否作用在号差改变时被解析延拓的那个时空坐标上。这一过程颇为微妙；众所周知，在量子场论中 $\sT$ 是反线性的，而 $\sR$ 是线性的。与此相关，$\sC$（电荷共轭）会进入延拓，因此欧几里得号差下的 $\sR$ 可以延拓为洛伦兹号差下的 $\sCT$，而 $\sCR$ 可以延拓为 $\sT$。若 $\sR$（或 $\sCR$）在费米子上平方为 $+1$，则 $\sCT$（或 $\sT$）平方为 $-1$，反之亦然。

为具体描述群 $\Pin^+(n)$，我们在 $\R^n$ 上工作，考察使一个坐标变号的反射如何作用在旋量上。令 $w$ 为单位向量，考虑反射 $x^i \to x^i -2w^i w\cdot x$。置 $\g\cdot w=\g_i w^i$，则该反射作用在旋量上为
\be\label{bel}\psi\to \g\cdot w\, \psi.\ee
这些变换平方为 $+1$，把它们与由 $\Sigma_{ab}=\frac{1}{4}[\g_a,\g_b]$ 生成的 $\Spin(n)$ 变换合并，就得到 $\Pin^+(n)$。描述 $\Pin^-(n)$ 的一种方式很简单：让给定的反射作用为 $\psi\to \i \g\cdot w\psi$，该操作平方为 $-1$。

为在 $W$ 上定义 $\pin^+$ 结构，需把切丛的过渡函数 $M_{\x\y}$ 从 $\O(n)$ 提升为取值于 $\Pin^+(n)$ 的函数 $\h M_{\x\y}$。在维数 $\geq 2$ 时，这一提升可能存在障碍（例如 $\RP^2$ 不容许 $\pin^+$ 结构）。若 $W$ 确实容许 $\pin^+$ 结构，该结构也未必唯一；与自旋结构一样，$\pin^+$ 结构之间相差一个 $H^1(W,\Z_2)$ 中元素的扭转。带有 $\pin^+$ 结构选取的流形称为 $\pin^+$ 流形。

若 $W$ 是 $\pin^+$ 流形，则过渡函数 $\h M_{\x\y}$（如前理解为 $2^{[n/2]}\times 2^{[n/2]}$ 矩阵）定义了一个向量丛 $\P\to W$。我们可以定义作为该丛截面的旋量场 $\psi$，狄拉克算子 $\D=\i\slashed{D}$ 可作用于 $\psi$。但它的某些性质与可定向情形不同，正如正文中所见、并将在附录 \ref{do} 中进一步讨论的那样。

若 $\P\to W$ 是一个 $\pin^+$ 结构，则与自旋情形一样，任何其他 $\pin^+$ 结构都是 $\P\otimes \rho$，其中 $\rho\to W$ 是某个实线丛。然而当 $W$ 不可定向时，会出现新的要素。此时存在一个典范的实线丛 $\veps\to W$，即定向丛：$\veps$ 沿闭回路 $\ell\subset W$ 的和乐为 $-1$，当且仅当绕 $\ell$ 一周时 $W$ 的定向被反转，否则为 $+1$。$\P'=\P\otimes \veps$ 是 $W$ 上的一个新 $\pin^+$ 结构，我们称之为与 $\P$ 互补的 $\pin^+$ 结构。具体地说，$\P'$ 的过渡函数 $\h M'_{\x\y}$ 由 $\P$ 的过渡函数 $\h M_{\x\y}$ 得到：只要 $\h M_{\x\y}$ 位于 $\Pin^+(n)$ 的定向反转分支中，就将其符号反转。

% [Tutor Note: 不可定向时定向丛 $\veps$ 提供一个新的不变量：$\P'=\P\otimes\veps$ 与 $\P$ 都是 $\pin^+$ 结构，二者在定向反转的过渡函数上相差一个符号。若 $W$ 可定向，则 $\veps$ 平凡、$\P'=\P$，故“互补结构”只是不可定向情形才有的新对象。它正是奇维时狄拉克算子把 $\P$ 与 $\P'$ 互换的原因，也解释了为何奇维 $\pin^+$ 流形上不存在作用在单一 $\P$ 上的自伴狄拉克算子。]

上面关于 $\pin^+$ 的一切说法对 $\pin^-$ 都有直接的对应，只是 $\pin^-$ 结构的障碍从 4 维才开始出现。

\subsection{狄拉克算子（Dirac Operators）}\label{do}

考虑 $\R^n$ 中的一个反射：
\be\label{porz}\sR\psi(x_1,x_2,\dots,x_n)=\g_1\psi(-x_1,x_2,\dots,x_n).\ee
不难看出，这样的反射与狄拉克算子
\be\label{orzo} \D=\i\slashed{D}=\i\sum_{j=1}^n\g^j D_j. \ee
反对易。

% [Tutor Note: 原文“不难看出”省略的验证是：$\g_1$ 与 $\g_j$（$j\ne1$）反对易、与 $\g_1$ 对易；同时 $D_1=\partial_1$ 在 $x_1\to-x_1$ 下变号，而 $D_j$（$j\ne1$）不变。逐项相乘后整体恰多出一个负号，故 $\sR\D+\D\sR=0$。一般地，任何定向反转等距都使 $\slashed{D}$ 变号，这是后文用 $\D'=\bg\slashed{D}$ 构造可下降算子、以及 $\eta$ 在定向反转下变号的前提。]

更一般地，在自旋流形 $\h W$ 上，任何定向反转的对称性都与 $\D$ 反对易。设 $\h W$ 具有一个自由作用且满足 $\tau^2=1$ 的定向反转对称性 $\tau$，则可从 $\h W$ 过渡到其商 $W=\h W/\Z_2$。若此外 $\tau$ 还是旋量丛 $\S\to \h W$ 的对称性且仍满足 $\tau^2=1$，则 $W$ 被赋予两个 $\pin^+$ 结构 $\P$ 与 $\P'$，它们可如下定义：$\P\to W$ 的截面 $\psi$ 是 $\S\to \h W$ 的满足
\be\label{milox} \tau\h\psi=\h\psi, \ee
的截面 $\h\psi$；而 $\P'\to W$ 的截面 $\psi'$ 是 $\S\to \h W$ 的满足
\be\label{ilox}\tau\h\psi=-\h\psi.\ee
的截面 $\h\psi$。由于 $\tau$ 作用在 $\S$ 上的符号取决于任意选取，$\P$ 与 $\P'$ 之间的关系是完全对称的。

取任意点 $p\in \h W$ 以及从 $p$ 到 $\tau(p)$ 的路径 $\h\ell$。对 $\tau$ 取商后，$\h\ell$ 投影为闭回路 $\ell\subset W$，绕该回路一周时 $W$ 的定向被反转。式（\ref{milox}）与式（\ref{ilox}）中的相对负号意味着，$\pin^+$ 丛 $\P$ 与 $\P'$ 绕 $\ell$ 的单值性相差一个符号。因此 $\P$ 与 $\P'$ 正是附录 \ref{unor} 意义下互补的 $\pin^+$ 丛。

每个 $\pin^+$ 流形 $W$ 都可由这一构造得到，此时 $\h W$ 是 $W$ 的定向双重覆盖。

现在设 $\psi$ 是 $W$ 上取值于某个特定 $\pin^+$ 结构 $\P$ 的费米子场。由上述讨论清楚可见，通常的狄拉克算子 $\D=\i\sum_k\g^k D_k$ 不能作为作用在 $\psi$ 上的自伴算子来定义。在覆盖空间 $\h W$ 上，$\D$ 与 $\tau$ 反对易，因此它把 $\P$ 的截面映为 $\P'$ 的截面，反之亦然。

对奇数维 $n$，这一情形无法补救：不存在作用在取值于 $\P$ 的场 $\psi$ 上的自伴狄拉克算子。其结果是，相应的路径积分 $Z_\psi$ 可以定义为某个复值反对称双线性型的普法夫量（Pfaffian），但不能定义为自伴算子的行列式。因此它是复值的，且一般而言会遭遇涉及模 1 复数的整体反常，而不只是实数 $\pm 1$。这一问题在正文第 \ref{anombo} 节中已有讨论。

% [Tutor Note: 奇维 $\pin^+$ 流形上没有自伴 $\D$，$Z_\psi$ 只能定义为复值反对称双线性型的 Pfaffian：Pfaffian 的相位本身并无绝对意义，只能相对于某个参考结构定义，因此 $Z_\psi$ 会带有模 1 的整体反常。这与正文讨论的 3 维自旋拓扑量子场论（配分函数如 $(-1)^\zeta$，或 2 维时的 $e^{-\i\pi\eta/2}$）以及 Dai--Freed 定理的配对机制直接相关：单独定义体拓扑项或边界费米子配分函数都可能不合法，而二者之积才有良好定义。]

对偶数维 $n$，情况不同：总存在作用在取值于任意给定 $\pin^+$ 结构 $\P$ 的费米子场上的自伴狄拉克算子。正文中多处用到这一点，这里给出细节。首先从带自旋结构 $\S$ 的可定向流形 $\h W$ 出发。由于 $\h W$ 可定向，可以定义作用在 $\S$ 上的手征算子 $\bg$（并导致分解 $\S=\S_+\oplus \S_-$，其中 $\S_+$ 与 $\S_-$ 是正、负手征旋量丛）。假定我们要求 $\bg^2=1$，则 $\bg$ 的恰当定义稍依赖于 $n$ 模 4 同余于 0 还是 2。例如在 2 维（以及任意 $k$ 的 $4k+2$ 维），定义
\be\label{toroz}\bg=\frac{\i}{2!}\epsilon^{ij}\g_i\g_j. \ee
在 4 维（以及任意 $k$ 的 $4k$ 维），略去因子 $\i$ 而定义
\be\label{orz}\bg=\frac{1}{4!}\epsilon^{ijkl}\g_i\g_j\g_k\g_l. \ee

% [Tutor Note: $\bg$ 的定义是否带 $\i$ 因子取决于 $n\bmod4$：$n\equiv2\pmod4$ 时取 $\frac{\i}{2!}\epsilon^{ij}\g_i\g_j$，$n\equiv0\pmod4$ 时取 $\frac{1}{4!}\epsilon^{ijkl}\g_i\g_j\g_k\g_l$。两种写法都是为了让 $\bg^2=1$、本征值为 $\pm1$，从而手征投影算子 $\frac{1\pm\bg}{2}$ 合法。若不按维度区分而硬套同一公式，会得到 $\bg^2=-1$ 或本征值为 $\pm\i$ 的“手征算子”，投影算子随之失效；这是跨维度计算时最常见的符号错误之一。]

在偶数维可定向流形 $\h W$ 上，我们可以不用通常的狄拉克算子 $\D=\i\slashed{D}$，而使用另一个狄拉克算子
\be\label{elf} \D'=\bg \slashed{D}. \ee
$\D'$ 与 $\D$ 一样是自伴的，且二者谱相同，因为它们是共轭的。事实上，若 $U=(1-\i\bg)/\sqrt 2$，则
\be\label{welf}\D'=U\D U^{-1}.\ee
这依赖于 $\bg$ 与 $\Spin(n)$ 在旋量表示上的作用对易，从而与隐藏在 $\D$ 定义中的 $\Spin(n)$ 联络对易。

% [Tutor Note: $\D'=\bg\slashed{D}$ 与 $\D$ 通过 $U=(1-\i\bg)/\sqrt2$ 共轭，故谱相同；$\bg$ 与 $\Spin(n)$ 对易保证 $\D'$ 仍由同一协变导数构成。$U=e^{-\i(\pi/4)\bg}$ 是手征相位的 $\pi/4$ 旋转，其幺正性依赖 $\bg^2=1$。由此也可反推出：只有在偶维、且按上一注释正确选取 $\bg$ 的 $\i$ 因子时，$\D'$ 才是自伴算子。]

因此在 $\h W$ 上，使用 $\D$ 或 $\D'$ 皆可。然而若要过渡到不可定向流形 $W=\h W/\Z_2$（其中 $\Z_2$ 由定向反转等距 $\tau$ 生成），则最好使用 $\D'$。事实上，由于 $\bg$ 与 $\D$ 在 $\tau$ 下均为奇，$\D'$ 为偶，因而可下降为作用在 $W$ 的 $\pin^+$ 丛 $\P$ 上的一个完全合适的自伴算子。

在正文中，对于偶数维 $\Pin^+$ 流形 $W$ 的自伴狄拉克算子，我们有时简单地写 $\D$ 而不写 $\D'$。这样做的理由是：在 $\h W$ 上，所谓 $\D$ 或 $\D'$ 实际上取决于克利福德代数表示的选取。我们可以定义新的 $\\gamma$ 矩阵
\be\label{werm} \g_i'=U\g_iU^{-1}   \ee
具体地
\be\label{zerm}\g_i'=-\i\bg\g_i=\begin{cases}-\eps_{ij}\g^j & {\mathrm{if}} ~n=2\cr
                                                                       \frac{\i}{3!}\epsilon_{ijkl}\g^j\g^k\g^l& {\mathrm{if}} ~n=4 .\end{cases}\ee
用新 $\\gamma$ 矩阵表示，$\D'$ 取标准形式
\be\label{termo}\D'=\i\sum_k\g'{}^kD_k . \ee
（协变导数 $D_k$ 在旧、新 $\\gamma$ 矩阵下形式相同。）

% [Tutor Note: 原文“in fact up on $\h W$”疑为“upon $\h W$”之排印笔误（源第 2706 行）。另一个易错点是：$\g_i'=U\g_iU^{-1}=-\i\bg\g_i$ 是幺正共轭，故仍厄米，不要误以为 $-\i\bg$ 因子破坏了厄米性；真正的代价是反射变换律变为式（\ref{yrf}）；式（\ref{zerm}）中 $n=2$ 与 $n=4$ 两行的形式差异，正是 $\bg$ 是否带 $\i$ 的直接后果，跨维度套用会出差号。]

然而，若使用新的 $\\gamma$ 矩阵，则反射下的变换律（\ref{bel}）（从而 $\Pin^+(n)$ 的任何定向反转元素下的变换律）会被修改。它变为
\be\label{yrf}\psi\to \i\bg'\g'\cdot w\,\psi. \ee
（这里 $\bg'$ 是由新 $\\gamma$ 矩阵构造的手征算子；它与 $\bg$ 相同。）事实上，这是 4 维时空中粒子物理里最常见的选择。\footnote{例如可参看 \cite{BD} 的式（2.33），不过该式对狄拉克费米子写入了任意相位，而且给出的是费米子在宇称（三个空间坐标全部反射，而非只反射一个）下的变换公式。}我们之所以反而从（\ref{bel}）出发，除了它在数学上相当标准之外，还因为它在许多目的下给出一个在一切维数中一致适用的形式体系。很难找到对所有目的都方便的费米子约定。

% [Tutor Note: 式（\ref{yrf}）所谓“变换律被修改”只是表象：把 $\g_i'=U\g_iU^{-1}$、$\bg'=\bg$ 代入可见 $\i\bg'\g'\cdot w$ 与旧基下的 $\g\cdot w$ 表示同一变换，改换的只是旋量分量的基。物理上真正改变的是“以哪一组 Clifford 表示为基本表示”，文献 \cite{BD} 式（2.33）采用的正是粒子物理惯用的相位；阅读时不要把“modified”误解为对称性本身被改变，否则会多算一个并不存在的符号。]

\section{低维例子（Examples In Low Dimension）}\label{examples}

\begin{keytakeaways}
\begin{itemize}[leftmargin=1.6em,itemsep=0.2em,topsep=0.2em]
  \item \textbf{物理图像}：低维例子把一般的 pin 结构理论落到具体维度。2 维旋量表示是实的，3 维是赝实的，4 维与 5 维仍是赝实；反线性操作 T 的平方以及是否存在整体手征算子，决定狄拉克谱如何成对，从而决定马约拉纳费米子配分函数的相位与整体反常。
  \item \textbf{数学工具}：借助 SU(2) 的赝实性、不变量张量、pin+ 与 pin- 两支的不可约表示、Kramers 型反线性算符，以及 APS 指标定理对 $\\eta$ 不变量的计算。
  \item \textbf{核心结论}：2 维有平方为 +1 的反线性操作 T（仅可定向时还有平方为 -1 的 T'）；3 维的 R 与 R' 是互为复共轭的复表示，T 在不可定向流形上把两个互补结构互换，故模 2 指标可非零；4 维只有一个自伴狄拉克算子且 T 平方为 -1；5 维指标在闭流形上为零、带 APS 边界时为偶数，构成 4 维 $\\eta$ 不变量讨论的基础。
\end{itemize}
\end{keytakeaways}

我们将更仔细地考察本文中重要的低维例子。

\subsection{二维（Two Dimensions）}\label{twod}

为描述 $\R^2$ 上的旋量，只需两个 $\\gamma$ 矩阵，可以取实 $2\times 2$ $\\gamma$ 矩阵 $\g_1=\sigma_1$，$\g_2=\sigma_3$。群 $\Spin(2)$ 由 $\Sigma=\frac{1}{2}\g_1\g_2$ 生成。抽象地说，$\Spin(2)$ 与 $\SO(2)$ 同构，旋量构成其基本的实二维表示。但由于 $\Sigma$ 的本征值为 $\pm \i/2$（而非 $\pm \i$），$\Spin(2)$ 是作用在 $\R^2$ 上作旋转的原始 $\SO(2)$ 的双重覆盖。

为得到 $\Pin^+(2)$，我们纳入反射 $x^i\to x^i-2w^i w\cdot x$，其作用为
\be\label{pes} \psi\to \gamma\cdot w\, \psi. \ee
这给出 $\Pin^+(2)$ 的一个实二维表示。群 $\Pin^+(2)$ 抽象地同构于 $\O(2)$，尽管它是作用在 $\R^2$ 上作旋转与反射的通常 $\O(2)$ 的双重覆盖。

天真地看，人们或许以为可以定义 $\Pin^+(2)$ 的第二个实二维表示，其中反射作用为
\be\label{nes}\psi\to -\gamma\cdot w\psi. \ee
很快我们会看到，类似的做法在 3 维确实可行，但在 2 维，这两个表示通过 $\psi\to \g_1\g_2\psi$ 共轭，后者代表 $\R^2$ 的一个 $\pi$ 旋转。

由于 $\Pin^+(2)$ 的旋量表示是实的，存在反线性操作
\be\label{wes}\T\psi =\psi^*,~~\T^2=1 \ee
它与厄米的狄拉克算子 $\D'$ 反对易，并保证在任何 2 维流形（可定向或否）上谱在 $\lambda\leftrightarrow -\lambda$ 下对称。

在{\it 可定向} 2 维流形上，我们还可以定义反线性操作
\be\label{res}\T'\psi=\bg \psi^*,~~~(\T')^2=-1 \ee
它与 $\D$（或 $\D'$）{\it 对易}，并保证狄拉克本征值都有偶重数。这里 $\bg=\i\g_1\g_2=\sigma_2$ 是手征算子。

% [Tutor Note: 2 维的关键在于旋量表示是实表示，故 $\T\psi=\psi^*$ 是反线性且平方为 $+1$ 的操作，使任意（可定向或不可定向）2 维流形上谱关于 $\lambda\leftrightarrow-\lambda$ 对称；可定向时另有 $\T'=\bg\psi^*$ 平方为 $-1$，给出偶重数。不可定向时 $\bg$ 没有整体定义（它依赖 $\S_+\oplus\S_-$ 的手征分解，而不可定向时该分解不存在），故 $\T'$ 不存在------这是理解 2 维 $\pin^-$ 分类（$\eta=\pm1/2$ 使 $e^{-\i\pi\eta/2}$ 为 8 次单位根）的几何出发点。]

两种结构之所以都存在，是因为若限制到 $\Spin(2)$，二维旋量表示可约；它分解为正、负手征的一维表示，二者互为复共轭。相应地，狄拉克算子与 $\bg$ 反对易，我们可以把它纳入 $\T'$ 的定义中。在不可定向 2 维流形上，$\T$ 的定义仍有意义，但 $\T'$ 不再能定义，因为 $\bg$ 没有整体定义。

\subsection{三维（Three Dimensions）}\label{wwod}

现在转到 3 维。对于 $2\times 2$ $\\gamma$ 矩阵，可以取 Pauli $\sigma$ 矩阵
\be\label{zob} \g_a=\sigma_a,~~~a=1,\dots,3.\ee
作用在旋量上的 $\SO(3)$ 的双重覆盖是 $\Spin(3)=\SU(2)$。旋量按 $\SU(2)$ 的自旋 1/2 表示变换。该表示是赝实的，因为 $\SU(2)$ 的平凡表示在张量积 $1/2\otimes 1/2$ 中{\it 反对称地}出现。具体地说，$\SU(2)$ 自旋 1/2 表示上的不变反对称张量由 $2\times 2$ Levi-Civita 符号 $\veps_{\alpha \beta}$，$\alpha,\beta=1,2$，给出。若 $\psi$ 是取值于该表示的费米子场，则相应的“质量项”（常写作 $\bar\psi\psi$）为
\be\label{mast}(\psi,\psi)=\veps_{\alpha\beta}\psi^\alpha\psi^\beta. \ee

群 $\Pin^+(3)$ 由在 $\Spin(3)=\SU(2)$ 上添加关于 $\R^3$ 中单位向量 $w$ 的反射矩阵 $\gamma\cdot w$ 得到。这样的反射矩阵是行列式为 $-1$ 的 $2\times 2$ 酉矩阵。因此 $\Pin^+(3)$ 是 $\U(2)$ 中行列式为 $\pm 1$ 的矩阵构成的子群：
\be\label{pinthree}\Pin^+(3)=\{g\in \U(2)|\det g=\pm 1\}. \ee
与 $\Spin(3)$ 不同，$\Pin^+(3)$ 有两个不等价的二维不可约表示 $R$ 与 $R'$。若 $\psi$ 与 $\psi'$ 分别取值于 $R$ 与 $R'$，则 $g\in \Pin^+(3)$ 对它们的作用为
\be\label{inth}\psi \to g\psi,~~  \psi'\to g(\det g)\,\,\psi'. \ee
因此差别恰恰在于：空间反射（或任何定向反转变换）作用在 $\psi$ 与 $\psi'$ 上时符号相反。表示 $R$ 与 $R'$ 都是复表示，即它们不容许不变双线性型。事实上，质量项 $(\psi,\psi)=\veps_{\alpha\beta}\psi^\alpha\psi^\beta$ 与 $(\psi',\psi')=\veps_{\alpha\beta}\psi^{'\alpha}\psi^{'\beta}$ 在反射下都是奇的。在 $\psi\to g\psi$，$\psi'\to g\det g\,\,\psi'$ 下，有
\be\label{zomb}(\psi,\psi)\to \det g\,\,(\psi,\psi),~~(\psi',\psi')\to \det g\,\,(\psi',\psi'). \ee
但联系两个在 $\sR$（或 $\sT$）下以相反方式变换的费米子 $\psi$ 与 $\psi'$ 的“非对角”质量项
\be\label{minth} (\psi,\psi')=\veps_{\alpha\beta}\psi^\alpha\psi^{'\beta} \ee
是不变的。

由于紧致群（如 $\Pin^+(3)$）的表示与其复共轭之间总存在不变配对，配对 $(\psi,\psi')$ 的不变性意味着表示 $R$ 同构于 $R'$ 的复共轭。事实上可以证明，作用在表示 $R$ 上的矩阵 $g$ 的复共轭，可通过一个幺正变换共轭为作用在表示 $R'$ 上的相应矩阵 $g\det g$：
\be\label{winth}\sigma_y \,g^*\sigma_y^{-1}=g\det g. \ee
为看出这一点，可以分别考虑 $g=\gamma\cdot w$ 为反射以及 $g\in \Spin(3)$（此时 $\det g=1$）两种情形。式（\ref{winth}）等价于 $\T g\T^{-1}=g\det g$，其中 $\T$ 是我们稍后引入的反线性操作。

% [Tutor Note: 原文“one can show”省略了 $\sigma_y g^*\sigma_y^{-1}=g\det g$ 的验证：对 $g=\gamma\cdot w$（$\det g=-1$）用 $\sigma_y\sigma_i^*\sigma_y^{-1}=-\sigma_i$ 可得结果为 $-g=g\det g$；对 $g\in\SU(2)$ 则用 $\SU(2)$ 的赝实性（即 $\sigma_yg^*\sigma_y^{-1}=g$）。这一恒等式说明 $R$ 与 $R'$ 互为复共轭，也说明非对角质量项 $(\psi,\psi')$ 为何在反射下不变，从而解释了 3 维边界费米子为何只能成对出现。]

在{\it 可定向} 3 维流形上，$\Spin(3)\cong \SU(2)$ 的旋量表示容许不变反对称张量 $\veps_{\alpha\beta}$，如上所述。因此可以定义反线性操作
\be\label{tork} (\T\psi)^\alpha=\veps^{\alpha\beta}\psi^*_\beta,\ee
其中 $\psi^*_\beta$ 是 $\psi^\beta$ 的复共轭。它满足 $\T^2=-1$，且与狄拉克算子 $\D$ 对易。于是有某种 Kramers 加倍：$\D$ 的本征值都有偶重数。

我们给出了相当内蕴的 $\T$ 定义，但一些读者可能还想看局部坐标下的描述。在选取标架与“平直空间” $\\gamma$ 矩阵 $\g_a=\sigma_a$、并把复共轭记为 $\K$ 之后，可以定义
\be\label{zork} \T=\K \sigma_y,~~\T^2=-1.\ee
该操作可用任意标架定义；并且由于它实际上是 $\SU(2)$ 不变的（尽管写法未使这一点显然），$\T$ 的定义不依赖于标架的选取，只要我们只考虑在时空上确定同一定向的标架（即只允许标架之间作保持定向的变换）。

% [Tutor Note: 原文源第 2819 行“betwen”为“between”之排印笔误，不影响文意。另需注意：$\T=\K\sigma_y$ 的写法只在保持定向的标架变换下不变；若允许定向反转的标架变换，$\T$ 会因与反射元素反对易而改变符号，这正是后文 $\T$ 不能在不可定向流形上整体定义的原因。]

可定向 3 维流形上狄拉克算子谱有偶重数这一事实的一个特例是：该算子的模 2 指标为零。然而 $\T$ 与 $\\gamma$ 矩阵反对易，从而与 $\Pin^+(3)$ 的反射元素反对易。因此，在不可定向的 3 维流形 $W$ 上，$\T$ 不能定义为旋量场之间的变换。上升到定向双重覆盖 $\h W$（$W=\h W/\Z_2$）后，$\T$ 可以定义，但与 $\Z_2$ 的定向反转生成元 $\tau$ 反对易。于是若 $\h\psi$ 是 $\h W$ 上的旋量场且 $\tau\h\psi=\pm \h\psi,$ 则 $\tau(\T\h\psi)=\mp \T\h\psi$。回顾互补 $\pin^+$ 结构 $\P$ 与 $\P'$ 的定义，我们可以在 $W$ 上解释这一陈述。它意味着：尽管 $\T$ 不能定义为把 $\P$ 的截面映到自身、或把 $\P'$ 的截面映到自身的映射，但作为从 $\P$ 的截面到 $\P'$ 的截面（及反之）的映射是有意义的。特别地，这意味着 $\P$ 与 $\P'$ 具有相同的模 2 指标（不必为 0，正如我们从第 \ref{inter} 节所知）。事实上，$\T$ 把狄拉克算子 $\D:\P\to \P'$ 的零模变为其伴随狄拉克算子 $\D:\P'\to\P$ 的零模。

% [Tutor Note: 不可定向时 $\T$ 与 $\tau$ 反对易，故它把 $\P$ 的截面映到 $\P'$ 的截面；于是 $\P$ 与 $\P'$ 的模 2 指标相等。原文“同样的模 2 指标”常被误读为“指标为零”：可定向时确实为零，不可定向时一般不为零（如 $\KB\times S^1$ 上的例子）。这正是 3 维自旋拓扑量子场论中配分函数 $(-1)^\zeta$ 的来源，也是后文把 $\eta$ 的模 2 约化与体拓扑项联系起来的伏笔。]

\subsection{四维（Four Dimensions）}\label{fwod}

在 4 维需要 4 个 $\\gamma$ 矩阵。可以取 $4\times 4$ 矩阵，但不可能全为实。例如可取
\be\label{dorf}\g_1=\sigma_1\otimes 1,~~\g_2=\sigma_2\otimes 1,~~\g_3=\sigma_3\otimes \sigma_1, ~~\g_4=\sigma_3\otimes \sigma_3.\ee
在此选择下，$\g_2$ 为虚，其余为实。

群 $\Spin(4)$ 同构于两个 $\SU(2)$ 的乘积，我们写作 $\SU(2)_\ell\times \SU(2)_r$。$\Spin(4)$ 的作用与手征算子 $\bg=\g_1\g_2\g_3\g_4$ 对易，因此 $\Spin(4)$ 的四维旋量表示分裂为两个带 $\bg=\pm 1$ 的二维表示之和。它们就是 $\SU(2)_\ell\times \SU(2)_r$ 的表示 $(1/2,0)$ 与 $(0,1/2)$。

这两个表示都是赝实的，原因很简单：$\SU(2)$ 的二维表示是赝实的。因此 $\Spin(4)$ 的旋量是赝实的。若把取值于 $(1/2,0)$ 表示的场记为 $\psi^{\alpha'}$，取值于 $(0,1/2)$ 表示的场记为 $\psi^{\alpha''}$，则对 $(1/2,0)$ 与 $(0,1/2)$ 表示分别有不变反对称张量 $\veps_{\alpha'\beta'}$ 与 $\veps_{\alpha''\beta''}$。在 $\SU(2)_\ell$ 或 $\SU(2)_r$ 对称性下，它们各自被确定到只差一个常数倍数。

为从 $\Spin(4)$ 过渡到 $\Pin^+(4)$，我们加入反射对称性 $x^i\to x^i-2w^i(w\cdot x)$。反射照常作用为
\be\label{zim}\psi\to \gamma\cdot w \,\psi.\ee
这种反射交换 $\SU(2)_\ell$ 与 $\SU(2)_r$，与 $\bg$ 反对易，并交换 $(1/2,0)$ 与 $(0,1/2)$。因此 $\Spin(4)$ 的表示 $(1/2,0)\oplus (0,1/2)$ 作为 $\Pin^+(4)$ 的表示是不可约的。$\Pin^+(4)$ 的这种旋量表示在同构意义下只有一个，因为式（\ref{zim}）中的一个符号可以通过与 $\bg$ 共轭而消去。

$\Spin(4)$ 不确定 $\veps_{\alpha''\beta''}$ 与 $\veps_{\alpha'\beta'}$ 之间的相对归一化。然而，交换 $\SU(2)_\ell$ 与 $\SU(2)_r$ 的反射把 $\veps_{\alpha'\beta'}$ 的一个特定选取映为 $\veps_{\alpha''\beta''}$ 的一个选取。在此选取下，和 $\veps_{\alpha'\beta'}\oplus \eps_{\alpha''\beta''}$ 是 $\Pin^+(4)$ 的四维旋量表示上的不变反对称双线性型 $\veps_{\alpha\beta}$。（指标 $\alpha,\beta$ 现在可以是 $\alpha',\beta'$ 型或 $\alpha'',\beta''$ 型。）

% [Tutor Note: 原文此处的求和式写作 $\veps_{\alpha'\beta'}\oplus \eps_{\alpha''\beta''}$，两个符号本应同为 $\veps$（$\eps$ 与 $\veps$ 都定义为 $\epsilon$，故不影响公式）；真正要注意的是：只有借助交换左右 $\SU(2)$ 的反射固定相对归一化，才能把两个分别不变的反对称张量拼成 $\Pin^+(4)$ 旋量表示上的单一不变反对称双线性型。若忘记这一步，会误以为 $\Spin(4)$ 已唯一确定 $\veps_{\alpha\beta}$。]

这使我们能在 $\pin^+$ 4 维流形 $W$ 上定义旋量的反线性操作：
\be\label{orf}(\T\psi)^\alpha =\veps^{\alpha\beta}\psi^*_\beta. \ee
由于 $\T^2=-1$，这种流形上自伴狄拉克算子 $\D'$ 的本征值有偶重数。

读者或许更愿意看到用局部标架显式定义的 $\T$。采用（\ref{dorf}）中平直空间 $\\gamma$ 矩阵的表示，定义为
\be\label{norf}\T=\K\gamma_2\bg. \ee
它与 $\Pin^+(4)$ 对易，因此该定义实际上不依赖于标架的选取，且具有整体意义。

% [Tutor Note: 4 维中 $\Pin^+(4)$ 的旋量表示唯一（$(1/2,0)\oplus(0,1/2)$ 在反射下不可约）且为赝实；$\T=\K\g_2\bg$ 与 $\Pin^+(4)$ 对易，故可整体定义，且 $\T^2=-1$ 给出 Kramers 型偶重数。与低维对照：2 维是实表示（$\T^2=1$）、3 维赝实（$\T^2=-1$）、4 维仍赝实；而 4 维 $\eta$ 取值为 $\pm1/4$，使 $e^{-\i\pi\eta/2}$ 成为 16 次单位根，这是附录 C 讨论 4 维 $\eta$ 的主要动机。]

\subsection{五维（Five Dimensions）}\label{fvod}

直观上，本文并不需要了解 5 维的旋量与狄拉克算子。但在附录 \ref{feta} 中，我们想知道 APS 指标定理对 4 维 $\eta$ 不变量的论断，为此需要了解 5 维的一些事实。

% [Tutor Note: 原文源第 2879 行在“eqn. (\ref{dorf})”之后误作句点，再另起一行“along with ...”，按文意应为逗号；这是排印笔误，不涉及公式内容。]

在 5 维可以找到克利福德代数的 $4\times 4$ 表示。使用式（\ref{dorf}）中的 4 个 $\\gamma$ 矩阵，再加上 $\g_5=\sigma_3\otimes \sigma_2$。注意 5 个 $\\gamma$ 矩阵中有 2 个是虚的（$\g_2$ 与 $\g_5$），其余为实。

因此群 $\Spin(5)$ 有一个四维旋量表示。\footnote{该表示在洛伦兹号差下是赝实的，因此需要取它的两份才能构造合理的费米子理论。这对我们此处有限的目的并不真正重要。}群 $\Spin(5)$ 实际上同构于 $\Sp(4)$，其四维旋量表示就是 $\Sp(4)$ 的基本四维表示。$\Sp(4)$ 定义为 $\U(4)$ 中保持某个非退化反对称双线性型的子群，因此特别地 $\Sp(4)$ 的四维表示是赝实的，$\Spin(5)$ 的旋量是赝实的。

具体地，在 $\Spin(5)$ 的旋量上可以定义反线性操作
\be\label{zanti}\T=\K\g_2\g_5,~~~ \T^2=-1.\ee
因此，在 5 维自旋流形（暂时可定向）上，存在某种 Kramers 加倍，狄拉克算子的所有本征值都有偶重数。

至此一切都与 3 维的情形平行。现在转向不可定向情形。反射按通常方式纳入。反射 $x^i\to x^i-2w^i(w\cdot x)$ 作用为
\be\label{wanti}\psi\to \pm (\gamma\cdot w)\psi. \ee
与 3 维一样，符号是有意义的，并给出 $\Pin^+(5)$ 的两个不同四维表示 $R$ 与 $R'$。然而有一个重要差别。在 3 维，$R$ 与 $R'$ 是互为复共轭的复表示；而在 5 维，它们各自是赝实的。为看出这一点，只需注意：式（\ref{zanti}）中定义的反线性操作 $\T$ 与按式（\ref{wanti}）作用的反射对易。因此 $\T$ 可以分别定义在 $R$ 与 $R'$ 中，且这两个表示都是赝实的。在初等层面上，与 3 维的差别在于：5 维中为虚的 $\\gamma$ 矩阵数目是偶数。

% [Tutor Note: 5 维中反线性操作 $\T=\K\g_2\g_5$ 与反射 $\pm\gamma\cdot w$ 对易，故 $\T$ 在 $R$、$R'$ 内各自有定义，两支都赝实；其代数原因是 $\T$ 扮演复共轭时，虚 γ 矩阵（此处为 $\g_2,\g_5$，共 2 个）的数目为偶数，正好与实 γ 矩阵的符号翻转相消。3 维中虚 γ 矩阵数目为 1（奇数），同样的构造只能把 $R$ 映到 $R'$，故两支互为复共轭而不各自赝实。这个“虚 γ 矩阵数目的奇偶”是判断实、赝实、复表示的最快办法。]

如同任何奇数维，不存在作用在取值于某个特定 $\pin^+$ 结构的场上的自伴狄拉克算子；相反，$\pin^+$ 结构以互补对 $\P$、$\P'$ 的形式出现，并被狄拉克算子交换。我们可以把狄拉克算子的{\it 指标} $\I$ 定义为 $\D:\P\to\P'$ 的零模数减去伴随算子 $\D:\P'\to\P$ 的零模数。在无边的五维流形上，该指标为 0；但在带边流形上（并采用 APS 边界条件），它可以非零，正如我们将在附录 \ref{feta} 中学到的那样。然而，由于反线性对称性 $\T$ 保证 $\D$ 作用在 $\P$ 或 $\P'$ 上的零模数都是偶数，$\I$ 总是偶数。

% [Tutor Note: 原文把 $\I$ 说成“零模数之差”，在闭五维流形上它确实为零（$\T$ 把 $\ker\D$ 与 $\ker\D^\dagger$ 配对，或等价地用 Hodge 理论看出）。带边时真正的对象是 APS 指标：定理断言 $\I$ 等于 $-\eta/2$ 在两端之差（更一般地还须计入边界零模数 $h$），这正是 4 维 $\eta$ 不变量取 $\pm1/4$ 的几何来源。$\T$ 使 $\D$ 的零模成对出现，故 $\I$ 恒为偶数；这一偶性保证了 $e^{-\i\pi\eta/2}$ 的 16 次单位根结构。]
